\section{Rendicontazione delle ore}

Questa sezione riporta le fasi affrontate nella realizzazione del progetto, insieme alla rendicontazione delle ore previste e di quelle effettivamente impiegate.

Per massimizzare la separazione imposta dall'architettura MVC, si è deciso di sviluppare l'intero motore logico prima dell'apprendimento del framework Qt. Fanno eccezione le classi \texttt{QString} e \texttt{QDateTime}, dato che non sono inerenti all'interfaccia grafica.

Per quanto riguarda la GUI, si è scelto di procedere con un approccio ``sul campo'': invece di imparare tutte le classi principali del framework, si è preferito suddividere lo sviluppo in micro-fasi: ad ognuna corrisponde una fase di apprendimento immediatamente seguita da quella di sviluppo.

\begin{table}[h!]
\centering
\renewcommand{\arraystretch}{1.3}
\begin{tabular}{|l|c|c|}
\hline
\textbf{Attività} & \textbf{Ore Previste} & \textbf{Ore Effettive} \\
\hline
Progettazione del modello & 3 & 2 \\
\hline
Sviluppo e test del modello & 4 & 3 \\
\hline
Studio e progettazione dello scheletro della GUI & 6 & 7 \\
\hline
Studio e implementazione del sistema \textit{Signals-Slots} & 6 & 8 \\
\hline
Test di segnali e slots (solo GUI) & 3 & 2 \\
\hline
Sviluppo del controller e connessione delle parti & 5 & 6 \\
\hline
Test e debug generale & 3 & 6 \\
\hline
Implementazione della persistenza & 3 & 2 \\
\hline
Aggiunta funzionalità extra & 2 & 3 \\
\hline
Test e debug finale & 2 & 3 \\
\hline
Stesura della relazione & 3 & 3 \\
\hline\hline
\textbf{Totale}  & \textbf{40} & \textbf{45} \\
\hline
\end{tabular}
\end{table}

Lo sforamento del monte ore totale, per quanto non eccessivo (pari al 12,5\% del monte ore previsto),
è stato causato da una leggera difficoltà di apprendimento dei concetti più ostici del framework
Qt, come il sistema \textit{Signals-Slots}, e dalla fase di test e debug in seguito alla connessione
dei componenti dell'architettura MVC.